%% 
%% Copyright 2007-2025 Elsevier Ltd
%% 
%% This file is part of the 'Elsarticle Bundle'.
%% ---------------------------------------------
%% 
%% It may be distributed under the conditions of the LaTeX Project Public
%% License, either version 1.3 of this license or (at your option) any
%% later version.  The latest version of this license is in
%%    http://www.latex-project.org/lppl.txt
%% and version 1.3 or later is part of all distributions of LaTeX
%% version 1999/12/01 or later.
%% 
%% The list of all files belonging to the 'Elsarticle Bundle' is
%% given in the file `manifest.txt'.
%% 
%% Template article for Elsevier's document class `elsarticle'
%% with numbered style bibliographic references
%% SP 2008/03/01
%% $Id: elsarticle-template-num.tex 272 2025-01-09 17:36:26Z rishi $
%%
\documentclass[authoryear,preprint,12pt]{elsarticle}
%% Use the option review to obtain double line spacing
%% \documentclass[authoryear,preprint,review,12pt]{elsarticle}

%% Use the options 1p,twocolumn; 3p; 3p,twocolumn; 5p; or 5p,twocolumn
%% for a journal layout:
%% \documentclass[final,1p,times]{elsarticle}
%% \documentclass[final,1p,times,twocolumn]{elsarticle}
%% \documentclass[final,3p,times]{elsarticle}
%% \documentclass[final,3p,times,twocolumn]{elsarticle}
%% \documentclass[final,5p,times]{elsarticle}
%% \documentclass[final,5p,times,twocolumn]{elsarticle}

%% For including figures, graphicx.sty has been loaded in
%% elsarticle.cls. If you prefer to use the old commands
%% please give \usepackage{epsfig}

%% The amssymb package provides various useful mathematical symbols
\usepackage{amssymb}
%% The amsmath package provides various useful equation environments.
\usepackage{amsmath}
%% The amsthm package provides extended theorem environments
%% \usepackage{amsthm}

%% The lineno packages adds line numbers. Start line numbering with
%% \begin{linenumbers}, end it with \end{linenumbers}. Or switch it on
%% for the whole article with \linenumbers.
%% \usepackage{lineno}

\usepackage{natbib}

\journal{Nuclear Physics B}



\begin{document}

\begin{frontmatter}

%% Title, authors and addresses

%% use the tnoteref command within \title for footnotes;
%% use the tnotetext command for theassociated footnote;
%% use the fnref command within \author or \affiliation for footnotes;
%% use the fntext command for theassociated footnote;
%% use the corref command within \author for corresponding author footnotes;
%% use the cortext command for theassociated footnote;
%% use the ead command for the email address,
%% and the form \ead[url] for the home page:
%% \title{Title\tnoteref{label1}}
%% \tnotetext[label1]{}
%% \author{Name\corref{cor1}\fnref{label2}}
%% \ead{email address}
%% \ead[url]{home page}
%% \fntext[label2]{}
%% \cortext[cor1]{}
%% \affiliation{organization={},
%%             addressline={},
%%             city={},
%%             postcode={},
%%             state={},
%%             country={}}
%% \fntext[label3]{}

\title{Master’s Thesis Proposal: Fairness Across Earliness Levels in Outcome-Oriented Predictive Process Monitoring}

%% use optional labels to link authors explicitly to addresses:
%% \author[label1,label2]{}
%% \affiliation[label1]{organization={},
%%             addressline={},
%%             city={},
%%             postcode={},
%%             state={},
%%             country={}}
%%
%% \affiliation[label2]{organization={},
%%             addressline={},
%%             city={},
%%             postcode={},
%%             state={},
%%             country={}}

\author{Alina Müller} %% Author name

%% Author affiliation
\affiliation{
    organization={Humboldt University of Berlin},
    city={Berlin},
    country={Germany}
}

%% Abstract
% \begin{abstract}
% %% Text of abstract
% Abstract text.
% \end{abstract}

%%Graphical abstract
% \begin{graphicalabstract}
% %\includegraphics{grabs}
% \end{graphicalabstract}

%%Research highlights
% \begin{highlights}
% \item Research highlight 1
% \item Research highlight 2
% \end{highlights}

%% Keywords
% \begin{keyword}
% %% keywords here, in the form: keyword \sep keyword

% %% PACS codes here, in the form: \PACS code \sep code

% %% MSC codes here, in the form: \MSC code \sep code
% %% or \MSC[2008] code \sep code (2000 is the default)

% \end{keyword}

\end{frontmatter}

%% Add \usepackage{lineno} before \begin{document} and uncomment 
%% following line to enable line numbers
%% \linenumbers

%% main text
%%

%% Use \section commands to start a section
\section{Background and Motivation}

I am currently completing a master's degree in Business Information Systems, focusing primarily on Data Science and secondarily on Business Processes. Through courses such as Business Analytics, Data Science, Business Process Automation, Conceptual Modeling, Applied Econometrics, and Deep Learning for Text Analytics, I have developed an interest in combining data-driven methods with process-oriented research. In the seminars Applied Predictive Analytics and Information Systems, I worked on machine learning methods for causal discovery. These projects provided me with practical experience implementing, evaluating, and comparing different machine learning approaches, further strengthening my interest in empirical research.

This interest is closely connected to my current work, in which I regularly use Python to automate data processing and operational tasks. Over time, I have become interested not only in automating individual activities but also in understanding how entire business processes can be analyzed, monitored, and improved.

After completing my master's degree, I would like to continue working in business processes and process automation. I am particularly interested in roles that combine process analysis, data science, and intelligent automation. Predictive process monitoring is highly relevant to my future career goals because it combines process data, machine learning, and operational decision support.

I have not yet taken the Process Prediction and Machine Learning course, but I believe my background in machine learning, predictive analytics, empirical analysis, and Python provides a solid foundation for working on this topic. I am confident that I can independently acquire the additional process-prediction knowledge required for the thesis, and I see this as a valuable opportunity to expand my expertise. I am also interested in taking the course during the upcoming winter semester while working on the thesis because it would help me address any remaining knowledge gaps and connect the course content to my research.

The proposed thesis would enable me to expand my knowledge of predictive process monitoring and process mining, building on my existing experience with machine learning, empirical analysis, and Python. Therefore, the topic creates a strong connection between my academic background, my current work, and the field in which I would like to develop professionally after graduation.


\section{Research Problem}
Predictive Process Monitoring supports decisions in high impact domains such as healthcare and finance, where outcome-oriented models predict binary or categorical case outcomes before a case has finished. A growing research stream examines fairness-related concerns in such predictions, showing that models trained on historical event logs can encode and reproduce discriminatory patterns against groups defined by sensitive attributes \citep{Peeperkorn2025FairnessThroughIndependencePPM, LeoniFairnessInPPAnalyticsAdversarialLearning, KappelHumanintheLoopApproach}. \citet{Revoredo2025FairPM} systematize this research area through FairPM, a taxonomy comprising 16 bias types and 12 mitigating interventions for
process mining, adapted from the machine learning fairness literature.

A second, largely unaddressed dimension concerns case earliness. Outcome-oriented PPM models are conventionally evaluated not only on complete cases but on running, partially observed cases (prefixes) at varying stages of completion, a practice established as the standard evaluation methodology \citep{TemporalStabilityPPM}. This earliness dimension matters for fairness: predictions made early in a case rely on less process information and may depend more heavily on static case attributes that serve as proxies for sensitive groups \citep{SequenceEncodingsPPM}, and it has been shown that an algorithm fair on complete data may exhibit disproportionate bias on incomplete observations \citep{fairnessAndMissingData}. Yet no existing fairness-in-PPM study has systematically examined how the group-fairness gap, or the fairness-accuracy trade-off obtained through bias-mitigation techniques, evolves as a case progresses. \citet{Peeperkorn2025FairnessThroughIndependencePPM} use prefix-based data but do
not report fairness metrics as trajectories over prefix length, \citet{LeoniFairnessInPPAnalyticsAdversarialLearning, KappelHumanintheLoopApproach} likewise do not analyze fairness variation across prefix lengths. Moreover, evidence from outside PPM suggests that pre-processing and post-processing mitigations behave differently under varying data conditions \citet{SanchezMarques2026FairnessGlioma}, raising the question of whether such differences also manifest across earliness levels. Closing this gap connects two currently separate strands of PPM research — fairness-aware prediction \citep{Peeperkorn2025FairnessThroughIndependencePPM, LeoniFairnessInPPAnalyticsAdversarialLearning} and earliness-aware prediction \citep{TemporalStabilityPPM, TradeOffAccurancyVSEarliness} — that have so far been studied independently.


\section{Research Question}

How does the fairness of outcome predictions, and the trade-off between fairness and accuracy, change as a case progresses from early to late? And do different bias-mitigation methods behave differently depending on how early the prediction is made?

This overarching question is broken down into five research goals, each building on the previous:

Goal 1: Baseline prediction pipeline. Establish a working outcome-prediction pipeline consisting of prefix generation, leakage-safe train/validation/test splitting, and a first trained model. Evaluate predictive performance (AUC, F1) across earliness levels, without any fairness component yet.

Goal 2: Baseline fairness gap. Using the pipeline from Goal 1, measure the group-fairness gap between sensitive-attribute groups at each earliness level, with no mitigation applied. This addresses whether a measurable, systematic relationship exists between case earliness and the magnitude of group unfairness.

Goal 3: Comparison across model families. Repeat the analysis from Goal 2 using three base models representing distinct PPM paradigms (classical, recurrent, attention-based). Assess whether the earliness-fairness relationship observed in Goal 2 differs by model architecture.

Goal 4: Effect of bias mitigation. Apply fairness-mitigation techniques on top of the models from Goals 2 and 3, and re-measure the fairness gap at each earliness level. Assess whether the relative effectiveness of different mitigation approaches (e.g., pre-processing vs. post-processing) changes with prefix length, and in which direction.

Goal 5: Fairness-accuracy trade-off. For each model, mitigation, and earliness combination from Goal 4, characterize the trade-off between the fairness improvement achieved and the accuracy cost incurred. Assess whether this trade-off improves or worsens as case completeness increases.

Goals 1 and 2 establish the necessary empirical baseline. Goals 3 to 5 constitute the main analytical contribution of the thesis, connecting fairness-aware and earliness-aware predictive process monitoring.

\section{Proposed Approach}
The analysis will be conducted using three base outcome-prediction models, each representing a distinct paradigm in predictive process monitoring: XGBoost, representing the classical machine-learning approach; an LSTM, representing the recurrent neural network paradigm; and ProcessTransformer \citep{Bukhsh2021ProcessTransformerPB}, representing the attention-based paradigm. Using multiple architectures serves two purposes: it assesses whether fairness behavior differs across model types, and it ensures that the findings on the earliness-fairness relationship are not specific to a single architecture.

Each model will be trained separately at three milestone-based earliness levels. For each earliness level and model, predictive performance (AUC, F1) and group-fairness metrics (e.g., demographic parity difference, equalized odds difference) will be measured with respect to the sensitive attribute. This addresses Goals 1 and 2, and provides the baseline for Goals 4 and 5.

Bias-mitigation methods (Goals 4 and 5). The specific mitigation techniques to be applied and compared remain to be finalized. Candidates under consideration, drawn from the FairPM taxonomy \citep{Revoredo2025FairPM} and the wider fairness-in-ML literature, include pre-processing, post-processing, and in-processing approaches. The final selection will be based on compatibility with the selected model families and the specific bias patterns observed in the baseline results.

The methodological choices outlined above are provisional and will be refined once the data is accessible and the computational feasibility of training and evaluating all model configurations within the thesis timeframe can be assessed. I am open to suggestions and recommendations from the supervisor side and would appreciate any guidance.


\subsection{\textbf{Earliness levels - milestone based:}}

Early = first occurrence of Triage

Mid = first occurrence of Vital Sign Check

Late = first occurrence of Medicine Dispensation 

\textbf{Why:}

First Triage = the moment the patient is first formally assessed and given an acuity level. This is the earliest point real prioritization has happened at all.

First Vital Sign Check = the first time objective physiological data (not just a subjective triage judgment) exists on the patient. This is a meaningful shift from "assessment" to "measured evidence."

First Medicine Dispensation = the first time actual treatment begins, not just assessment. This marks the shift from evaluating the patient to acting on them.

1. First Triage = the patient walks in and a nurse looks at them and says "okay, how urgent is this?" That's the very first judgment call — how serious does this look.

2. First Vital Sign Check = someone actually measures things — blood pressure, heart rate, temperature. Now you have real numbers, not just a first impression.

3. First Medicine Dispensation = the patient actually gets medicine. This is the point where the hospital stops just looking at the patient and starts doing something to help them.

First occurrence = something a real system could actually use, right when it happens.

Last occurrence = something you can only know afterward, which isn't realistic for a live prediction.

\subsubsection{\textbf{Outcome variable: discharge decision}}
What about the 70\% "Unknown" discharge cases? These can't be used as labels at all, since you don't know the true outcome — they get dropped from model training/evaluation, not filed under either group.

\subsubsection{\textbf{Sensitive variables: insurance, language (optional secondary ) }}

\subsubsection{\textbf{Bias}}
1. Data Bias — the problem is already baked into your data, before any model touches it.

Example: if insurance status already correlates with different treatment historically, your data just reflects that unfairness — the model didn't invent it, it inherited it.

2. Algorithm Bias — the model itself introduces or amplifies unfairness while learning.

Example: a "one-size-fits-all" model that's accurate on average but performs worse for smaller groups, because it wasn't built to treat groups separately.

Pre-processing (fix the data before training): Reweighing, Sampling, Variable Blinding, Transformation, etc. — you touch the data, not the model.

In-processing (fix it during training): Adversarial Learning, Regularization, Bandits — you change how the model learns.

Post-processing (fix it after training): Calibration, Thresholding, Visualization — you adjust the model's output, leaving the model itself untouched.


Reweighing = pre-processing → touches your data

Calibration = post-processing → and this one is literally the same thing as "Equalized Odds," one of your two fairness metrics. So picking Calibration as your intervention and Equalized Odds as your metric is really "use the same idea twice" — which is a clean, easy-to-explain design choice.

Adversarial Learning = in-processing → the harder, architecture-specific one you're treating as a stretch goa

Reweighing (pre-processing)

Why this one specifically, not one of the other pre-processing options (Variable Blinding, Sampling, Transformation, etc.)?

It's the least invasive. It doesn't change your data values or remove anything — it just tells the model "pay more attention to this group's cases." The data itself stays real and untouched, which matters for a clinical dataset where you don't want to be inventing or deleting patient records.
It's model-agnostic. Since you're comparing three very different architectures (XGBoost, LSTM, ProcessTransformer), you need a mitigation method that means the same thing no matter which one you apply it to. Reweighing just adjusts training-instance importance — every model type can use that concept.
It's already proven in this exact context. Kate's own taxonomy cites da Silva et al. (2024) specifically testing Reweighing for predictive process monitoring — so you're not the first to try it here, you're extending a validated approach to a new angle (earliness).

The core idea, simply

Normally, when you train a model, every training example (every patient case) counts equally. Reweighing changes that: some cases count more, some count less, based on how rare or common their combination of sensitive-group + outcome is.

The concrete steps

1. Split your training data into four groups, based on two things: insurance status (insured / uninsured) and the actual outcome (home / facility). That gives you four buckets:

Insured + Home
Insured + Facility
Uninsured + Home
Uninsured + Facility

2. Count how big each bucket actually is, and compare that to what it "should" be if insurance and outcome were completely unrelated (statistically independent). For example, if uninsured patients are underrepresented in the "Home" bucket relative to what you'd expect by chance, that bucket needs to count for more.

3. Calculate a weight for each bucket using a simple formula:

weight = (expected proportion) / (actual proportion)

If a bucket is underrepresented, its weight is greater than 1 (each case in it counts extra). If a bucket is overrepresented, its weight is less than 1 (each case counts a bit less).

4. Assign that weight to every case in the bucket. So every "uninsured + Home" case, for example, gets the same weight — say, 1.4 instead of 1.0.

5. Pass those weights into your model's training step. This is the part that's genuinely easy to implement — most model libraries already support this natively:

XGBoost: has a built-in sample\_weight parameter you just pass in.
LSTM / ProcessTransformer (neural networks): most loss functions (e.g., cross-entropy) support per-example weighting directly — you multiply each example's loss by its weight before the model updates.

So practically, you're not changing your model's architecture at all — you're just adding one extra column (the weight) that your training function already knows how to use.

What changes across your three earliness levels

Here's the part that's actually relevant to your specific thesis: you'd recompute the weights separately at each earliness level (Triage / Vital Sign Check / Medicine Dispensation), not just once. Why? Because the composition of your data can shift at each cutoff — remember, some cases don't reach the later milestones at all, so the insured/uninsured, home/facility proportions among "cases that reached Medicine Dispensation" might look different from the full dataset. Recomputing weights per earliness level keeps the intervention honest to whatever data is actually available at that stage.

Other pre-processing options you're skipping

Variable Blinding — just remove insurance status from the training data entirely.
Why not: this defeats the purpose of your thesis. If insurance isn't in the data at all, you can't measure whether its influence changes across earliness levels — you'd have nothing to track. It's a legitimate fairness technique in general, but it's incompatible with a research question about how much a sensitive variable's influence shifts.

Sampling — split the data into subgroups (e.g., insured vs. uninsured) and train separate models or process representations per group.
Why not: this means building essentially separate pipelines per group per earliness level per architecture — the complexity multiplies fast, and Kate's own taxonomy notes it's only really been tried for process discovery so far, not predictive monitoring like yours. It's a bigger methodological leap with less precedent behind it.

Transformation — change the sensitive variable's values (or proxies) directly, e.g., mapping numeric attributes to preserve only order, not exact values.
Why not: it's really designed for continuous/numeric proxy variables (like the ZIP-code example in the paper). Insurance status is categorical, so this technique doesn't map naturally onto your case the way Reweighing does.

Relabeling and Perturbation — change the outcome labels themselves to balance fairness.
Why not: this is riskier in a clinical context specifically — you'd be artificially changing what "actually happened" to a patient's discharge outcome to make the numbers look fairer. That's a much harder thing to justify ethically and methodologically than just reweighting how much attention a case gets.

Causal Methods — identify statistical relationships between insurance and other variables first, to guide decisions.
Why not: this is more of a diagnostic/exploratory tool than a mitigation technique — genuinely useful (you're informally already doing a version of this by looking at the Andreswari paper's effect sizes), but it doesn't by itself fix anything. It's a step before choosing an intervention, not an intervention itself.

Calibration (post-processing)

It directly matches a metric you're already using. You're planning to report Equalized Odds as a fairness metric. Kate's taxonomy states plainly that Calibration is the intervention version of Equalized Odds — same underlying idea, just applied as a fix instead of just a measurement. That means your evaluation and your intervention are telling the same consistent story, instead of using one concept to measure and a totally different concept to fix.
It's also architecture-agnostic, for the same reason as Reweighing — it works on the model's output, after the model is already done, so it doesn't care what kind of model produced that output.
It's simple to explain and defend. You're not inventing a new fairness concept — you're just closing the loop on one you're already using.

Equalized odds asks a sharper question: among patients who actually needed to go home, does the model correctly send them home at the same rate for both groups? And among patients who actually needed a facility, does the model correctly flag that at the same rate?

Other post-processing options you're skipping

Thresholding — adjust the decision threshold per group (e.g., uninsured patients need a lower predicted-risk score to trigger the same outcome).
Why not: this is a real, legitimate alternative to Calibration, and worth naming as such — but Calibration wins because it maps directly onto Equalized Odds, a metric you're already reporting. Thresholding doesn't have that same built-in link to your evaluation metrics, so choosing it would mean explaining two separate concepts (a fairness metric and a differently-framed fix) instead of one consistent one.

Visualization — adjust how results are displayed to help users notice unfair patterns.
Why not: straightforwardly out of scope — this targets User Bias, and your thesis isn't studying how doctors interpret or act on model outputs, it's studying the model's own predictions. Not applicable to your research question at all.

Here's Calibration explained simply, then how you'd actually do it.

The core idea

Calibration is a post-processing fix — meaning your model is already fully trained, and you're not touching the training process at all. Instead, you take the model's predictions after the fact and adjust them so they're fair across groups.

Remember: Calibration in this taxonomy is the same thing as Equalized Odds — the goal is to make the model's error rates equal across insured and uninsured patients. Specifically: the rate of correctly predicting "home" among patients who actually went home, and the rate of incorrectly predicting "home" among patients who actually needed a facility — these should be roughly the same for both groups.

Why you need this at all

Your model outputs a probability — say, "78% chance this patient goes home" — not a hard yes/no. Somewhere, you turn that probability into a decision using a cutoff (e.g., "above 50% → predict home"). The trick behind Calibration is: you don't have to use the same cutoff for both groups.

If the model's errors are unbalanced — say, it's systematically too cautious about sending uninsured patients home — you can compensate by using a slightly different cutoff for that group specifically. Not because the patients are different, but because the model's mistakes are lopsided, and adjusting the cutoff corrects for that.

The concrete steps

1. Use a separate slice of your data — not the training set. You need a validation set the model hasn't seen, so you can measure its real-world error pattern honestly.

2. Get the model's predicted probabilities on that validation set, separately for insured and uninsured patients.

3. Measure the error rates for each group at a default cutoff (e.g., 0.5): how often does it correctly predict "home" for people who really went home, and how often does it wrongly predict "home" for people who really needed a facility.

4. If the two groups' error rates differ, search for a different cutoff for each group that brings the error rates closer together. This is usually done by trying a range of possible cutoffs (e.g., 0.3, 0.4, 0.5, 0.6…) for each group and picking the one that best equalizes the error rates between them — this is a small optimization step, not something you compute by hand.

5. Apply those two separate cutoffs going forward — insured patients get evaluated against their cutoff, uninsured patients against theirs, even though it's the exact same trained model underneath.

What's genuinely simple about this

You don't retrain anything. You don't touch your XGBoost/LSTM/ProcessTransformer models at all — this step happens entirely after you already have predictions. That also means it's the easiest of your three methods to apply consistently across all three architectures, since it only cares about the final probability output, not how the model produced it.

What changes per earliness level, for you specifically

Same as with Reweighing — you'd calibrate separately at each of your three milestones (Triage / Vital Sign Check / Medicine Dispensation), since the model's error pattern at an early prediction (less information) is likely very different from its error pattern later (more information). This actually ties directly into your core research question: does the cutoff adjustment needed to fix fairness get smaller or larger as more clinical information becomes available? That's basically Goal 5 in your proposal — the fairness-accuracy trade-off — measured concretely.


Adversarial Learning (in-processing, stretch goal)

It's the "proper" fairness technique in the literature — it directly trains the model to not be able to guess the sensitive attribute from its internal representations, which is a more thorough fix than just reweighting or recalibrating.
But it's architecture-specific. Adversarial learning needs a neural network with a competing "adversary" network attached — this is natural for your LSTM and ProcessTransformer, but genuinely awkward for XGBoost (which isn't trained by gradient descent the same way).
That's exactly why it's optional, not core. If you commit to it as a required method, comparing "in-processing effectiveness" across your three architectures becomes questionable — you'd partly be comparing "how well does adversarial learning work" vs. "does adversarial learning even mean the same thing here," which muddies your earliness-focused research question. Keeping it as a stretch goal means if time allows, you get a richer result; if not, your core two-method comparison (Reweighing vs. Calibration) is already clean and complete without it.

Other in-processing options you're skipping

Bandits — designed for fair online decision-making, learning fairness with little or no training data, in a live/sequential setting.
Why not: your setup isn't online/sequential decision-making — you have a large historical dataset to train on. Bandits solve a different problem (learning fairness on the fly with almost no data) that doesn't match your situation at all.

Regularization — add a fairness penalty term directly into the model's loss function.
Why not: this is actually a reasonable alternative to Adversarial Learning, and genuinely close in spirit. The reason to prefer Adversarial Learning if you do go in-processing at all is that it has a closer, more direct precedent in this exact research area (de Leoni et al., already in your related work) — but Regularization is worth mentioning to Kate as "the other in-processing option I considered," since it's arguably simpler to implement than adversarial learning if you do decide to pursue a stretch goal.

Reweighing (as in-processing) — the taxonomy notes reweighing can also be done as an in-processing step, where the model learns the weights itself during training.
Why not: you're already using Reweighing as pre-processing — doing it again as in-processing would be redundant rather than adding a genuinely different method to compare.


In-processing: Regularization (stretch goal)

Overall approach: if time allows after the core Reweighing/Calibration comparison (Goals 1-5) is complete, add Regularization as a third mitigation method — applied consistently across all three architectures, using the technique appropriate to each:

Model	Regularization implementation
XGBoost	FairGBM (LightGBM-based) — custom "constrained cross-entropy" objective enforcing equalized odds directly during training
LSTM	Standard regularization — add a fairness penalty term (e.g., based on equalized-odds gap) directly into the neural network's loss function
ProcessTransformer	Same approach as LSTM — fairness penalty term added to the loss function

\subsection{\textbf{Models}}



\section{Data}

The intended dataset is MIMICEL, a process-mining event log derived from the MIMIC-IV Emergency Department (MIMIC-IV ED) database. It contains genuine demographic case attributes rather than requiring artificial or proxy sensitive attributes. \citet{AndreswariMIMICEL} report that MIMICEL comprises 7,422,277 event rows across 413,893 cases, following a simple six-activity ED process (Enter ED, Triage, Vital Sign Check, Medicine Reconciliation, Medicine Dispensation, Discharge).

\textbf{Outcome and sensitive attribute.} Outcome and sensitive attribute. Based on the empirical analysis by \citet{AndreswariMIMICEL} on the same underlying data, the outcome variable will be the discharge decision (disposition), binarized into a clinically meaningful split (e.g., home discharge vs. further care or facility). The sensitive attribute will be insurance status, since \citet{AndreswariMIMICEL} found it to be associated with the largest and most consistent effect sizes among all tested demographic attributes (age, gender, race, language, insurance) for this outcome, including a very large effect at the lowest acuity level. This provides a substantially stronger empirical basis than attributes such as race or gender, which showed consistently negligible effects in their analysis.

\textbf{Limitations.} The percentage of "Unknown" values reported by Andreswari et al. for insurance status and for discharge decision is known individually (approximately 55\% and 70\% of cases, respectively). However, the overlap between these two subsets, i.e., how many cases have both a known insurance value and a known discharge decision, is not yet known and can only be determined once data access is obtained.

\textbf{Access note.} Access to MIMIC-IV and MIMICEL requires PhysioNet credentialing, including completion of a recognized human-subjects research training course and a signed data use agreement. This has not yet been completed. 


\bibliographystyle{apalike}
\bibliography{cas-refs}

\end{document}

\endinput
%%
%% End of file `elsarticle-template-num.tex'.
